% Chapter 1: Robust Bayesian Inference on Riemannian Submanifold
\chapter{黎曼子流形上的稳健贝叶斯推断}\label{ch:chapter1}

% ----------------------------------------------------------
\section{引言}\label{sec:ch1-intro}

统计推断的根本挑战之一在于：当数据的高维特征空间 $\mathbb{R}^D$
本质上由低维流形结构主导时，
经典方法往往失效。环境维度 $D$ 巨大（图像、文本、点云数据），
但有效自由度 $d \ll D$，数据点近乎均匀分布在某个未知低维黎曼流形 $\mathcal{M}$ 上。

在贝叶斯框架下，这一几何约束引出了一个深刻问题：
如何对取值于黎曼流形 $\mathcal{M}$ 的参数进行稳健的贝叶斯推断？
特别地，当数据生成机制偏离标准假设时——例如，
数据服从 Huber 污染模型（contamination model）——
后验分布是否仍能提供可靠的不确定性量化？

本论文（\citep{tang2023robust}）的核心贡献在于提出
Bayesian RPETEL 框架，并建立了一套完整的三重稳健性理论。

% ----------------------------------------------------------
\subsection{Huber 污染模型}\label{sec:ch1-huber}

设真实数据分布 $P$ 为纯净分布 $P_0$ 与污染分布 $Q$ 的混合：
\begin{equation}\label{eq:ch1-huber-model}
P = (1-\epsilon)\,P_0 + \epsilon\,Q,
\qquad 0 \le \epsilon < \frac{1}{2},
\end{equation}
其中 $\epsilon$ 为污染率（contamination rate）。
纯净分布 $P_0$ 关于勒贝格测度绝对连续，
其密度函数记为 $p_0$；污染分布 $Q$ 任意。
该模型刻画了"大部分数据来自正确模型，但存在少量对抗性或异常噪声"这一实际情形。

在传统欧氏空间设定下，Huber 损失与污染模型催生了
鲁棒似然（robust likelihood）和鲁棒估计的理论。
本文将此框架推广至黎曼流形 $\mathcal{M}$ 设定，
引入了流形约束下的欧氏风险最小化（Projected Euclidean Risk Minimization，
简称 PERM）以及其贝叶斯版本 RPETEL。

% ----------------------------------------------------------
\subsection{三重稳健性目标}\label{sec:ch1-triple-robust}

本文所追求的稳健性包含三个层面：

\begin{enumerate}
    \item \textbf{似然稳健性（likelihood robustness）}：
          即使似然函数未正确指定，后验分布仍能提供
          自动校准的不确定性量化。
    \item \textbf{风险稳健性（risk robustness）}：
          即使欧氏风险最小化器（Euclidean risk minimizer）
          不在流形 $\mathcal{M}$ 上，后验中心仍能接近
          流形上的风险最小化器。
    \item \textbf{先验稳健性（prior robustness）}：
          在正确设定（correct specification）下，
          后验分布与标准贝叶斯后验渐近等价，
          不因引入稳健机制而损失效率。
\end{enumerate}

这三重稳健性目标之间存在内在张力——
提高似然稳健性通常会牺牲效率，
而宽松的先验稳健性条件又可能削弱风险稳健性。
本文的核心结果表明，在一定正则性条件下，
Bayesian RPETEL 框架可以同时实现这三个目标。

% ----------------------------------------------------------
\section{RPETEL 框架}\label{sec:ch1-rpetel}

% ----------------------------------------------------------
\subsection{从 PERM 到 RPETEL}\label{sec:ch1-perm}

设损失函数 $\ell_x:\mathcal{M}\to\mathbb{R}_+$，风险函数定义为
\begin{equation}\label{eq:ch1-risk}
R(\theta) = \mathbb{E}_P[\ell_X(\theta)]
          = \int \ell_x(\theta)\,\mathrm{d}P(x).
\end{equation}
在欧氏空间中，风险最小化器（ERM）为
$\theta^\dagger = \arg\min_{\theta\in\mathbb{R}^D} R(\theta)$。
然而，当参数自然地取值于流形 $\mathcal{M}$ 时，
将 $\theta^\dagger$ 投影到 $\mathcal{M}$ 上是自然的想法。

\textbf{PERM}（Projected Euclidean Risk Minimization）的定义为：
\begin{equation}\label{eq:ch1-perm}
\theta^\ddagger = \arg\min_{\theta\in\mathcal{M}}\,R(\theta),
\qquad
\text{即 } \theta^\ddagger = \Pi_{\mathcal{M}}(\theta^\dagger),
\end{equation}
其中 $\Pi_{\mathcal{M}}(\cdot)$ 为到流形 $\mathcal{M}$ 的最近点投影。

\textbf{RPETEL}（Bayesian Robust Projected Euclidean Risk Minimization）
则将 PERM 嵌入贝叶斯框架：
设先验分布 $\Pi$ 在 $\mathcal{M}$ 上定义，
定义伪似然（pseudo-likelihood）如下：
对于观测数据 $x_1,\ldots,x_n$，定义加权负对数似然
\begin{equation}\label{eq:ch1-pseudolikelihood}
L_n(\theta) = \frac{1}{n}\sum_{i=1}^n \ell_{x_i}(\theta),
\qquad \theta\in\mathcal{M}.
\end{equation}
后验分布定义为
\begin{equation}\label{eq:ch1-posterior}
\Pi_n(\mathrm{d}\theta)
= \frac{\exp\!\big(-n L_n(\theta)\big)\,\Pi(\mathrm{d}\theta)}
       {\int_{\mathcal{M}}\exp\!\big(-n L_n(\theta')\big)\,\Pi(\mathrm{d}\theta')}.
\end{equation}

关键观察：在平方损失 $\ell_x(\theta)=\|x-\theta\|_2^2$ 下，
风险函数 $R(\theta)$ 的形式关于 $\theta$ 是二次的，
这使得后续的局部坐标分析和切空间投影变得 tractable。

% ----------------------------------------------------------
\subsection{平方损失下的风险函数}\label{sec:ch1-risk-square}

设损失函数 $\ell_x(\theta) = \|x-\theta\|_2^2$。
记 $m := \mathbb{E}_P[X] \in \mathbb{R}^D$ 为数据的均值向量，
则风险函数简化为
\begin{equation}\label{eq:ch1-risk-quadratic}
R(\theta) = \mathbb{E}_P[\|X-\theta\|_2^2]
          = \mathrm{tr}(\Sigma) + \|\theta-m\|_2^2,
\end{equation}
其中 $\Sigma = \mathrm{Cov}_P(X)$ 为协方差矩阵。
该分解揭示了一个核心事实：
$\theta^\dagger = m$，而 $\theta^\ddagger = \Pi_{\mathcal{M}}(m)$——
流形上的风险最小化器是均值向量在流形上的投影。

在后验分布 \cref{eq:ch1-posterior} 中，
指数项 $-n L_n(\theta)$ 对应于样本均值的指数型族，
这与欧氏空间中的经典结果形成类比。

% ----------------------------------------------------------
\section{黎曼流形 BvM 定理}\label{sec:ch1-bvm}

贝叶斯渐近理论的核心问题是：
当样本量 $n\to\infty$ 时，后验分布如何收缩到真实参数？
经典的 Bernstein-von Mises（ BvM）定理表明，
在正则条件下，后验分布渐近正态且集中于真参数附近，
中心为 $\sqrt{n}(\theta-\theta_0)$ 的渐近协方差由 Fisher 信息决定。

本文将这一经典结果推广到黎曼子流形 setting，
建立了流形上的 BvM 定理——这是本文最重要的理论贡献。

% ----------------------------------------------------------
\subsection{切空间投影与局部坐标}\label{sec:ch1-tangent}

设 $\mathcal{M}$ 为 $\mathbb{R}^D$ 中紧致光滑 $d$ 维黎曼子流形，
嵌入维度为 $D$。对任意点 $\theta\in\mathcal{M}$，
记其切空间（tangent space）为 $T_\theta\mathcal{M}\subset\mathbb{R}^D$，
维数为 $d$；记到切空间的正交投影算子为
$\Pi_\theta: \mathbb{R}^D \to T_\theta\mathcal{M}$。

给定后验分布 $\Pi_n$，定义其在切空间 $T_\theta\mathcal{M}$ 上的投影为
\begin{equation}\label{eq:ch1-posterior-projection}
\Pi_n^\theta(B) = \Pi_n\!\big(\{\theta'\in\mathcal{M}:
\Pi_\theta(\theta'-\theta)\in B\}\big),
\qquad B\subset T_\theta\mathcal{M}.
\end{equation}

局部坐标的构造依赖指数映射（exponential map）
$\exp_\theta: T_\theta\mathcal{M}\to\mathcal{M}$：
对任意切向量 $v\in T_\theta\mathcal{M}$，
$\exp_\theta(v)$ 为从 $\theta$ 出发沿测地线以速度 $\|v\|$ 行进得到的点。
在 $\theta$ 的小邻域内，指数映射为微分同胚，
因此可将 $\mathcal{M}$ 上的局部问题拉回到切空间 $T_\theta\mathcal{M}$ 上分析。

% ----------------------------------------------------------
\subsection{渐近正态性}\label{sec:ch1-asymptotic}

设真实参数为 $\theta_0\in\mathcal{M}$，
数据生成分布为 $P_{\theta_0}$（正确设定）。
以下定理给出了流形上 BvM 定理的精确表述。

\begin{Theorem}[黎曼流形 BvM 定理]\label[theorem]{thm:ch1-bvm}
假设以下正则性条件成立：
\begin{enumerate}
    \item[(A1)] 流形 $\mathcal{M}$ 为 $C^3$ 光滑紧致黎曼子流形；
    \item[(A2)] 损失函数 $\ell_x(\cdot)$ 在 $\mathcal{M}$ 上满足
          二次增长条件（quadratic growth condition）；
    \item[(A3)] 先验分布 $\Pi$ 在 $\theta_0$ 的邻域内具有正的光滑密度；
    \item[(A4)] $n\to\infty$ 时收缩条件 $\epsilon_n\to 0$ 成立，
          其中 $\epsilon_n$ 为某种熵量度。
\end{enumerate}
记 $\Pi_n^{\theta_0}$ 为后验分布在切空间 $T_{\theta_0}\mathcal{M}$ 上的投影
（\cref{eq:ch1-posterior-projection}）。
则
\begin{equation}\label{eq:ch1-bvm-limit}
\Pi_n^{\theta_0}
\;\Longrightarrow\;
N\!\big(0,\;n^{-1}\,G(\theta_0)^{-1}\big)
\quad\text{（在分布意义下），}
\end{equation}
其中 $G(\theta_0)$ 为黎曼度量（metric）在 $\theta_0$ 处的 Fisher 信息矩阵
在局部坐标下的表示。
\end{Theorem}

\cref{eq:ch1-bvm-limit} 的含义是：
将后验分布投影到 $\theta_0$ 处的切空间后，
该投影渐近等价于均值在 $\theta_0$、协方差为
$n^{-1}G(\theta_0)^{-1}$ 的正态分布。
这与欧氏空间 BvM 定理的形式完全类比，
唯一的差异在于协方差由黎曼 Fisher 信息矩阵 $G(\theta_0)$ 决定，
而非欧氏 Fisher 信息。

\textbf{证明思路}\footnote{详细证明见 \citep[Section 3]{tang2023robust}。
核心步骤分为三步：
（1）在局部坐标下将流形上的问题转化为切空间 $\mathbb{R}^d$ 上的问题；
（2）运用欧氏 BvM 定理的标准技术（van der Vaart \& Wellner 方法）证明切空间上的渐近正态性；
（3）通过指数映射 $\exp_{\theta_0}$ 的近似性质（在小邻域内接近线性映射）
将切空间结果拉回流形。其中步骤（1）需要用到流形的局部平坦化性质
和投影算子 $\Pi_\theta$ 的光滑性。}：

% ----------------------------------------------------------
\subsection{收敛速率}\label{sec:ch1-rate}

由 \cref{eq:ch1-bvm-limit}，后验分布在切空间投影上的集中速率由
$n^{-1/2}$ 主导。进一步的分析表明，
在额外正则性假设下，后验收缩速率（posterior contraction rate）为
\begin{equation}\label{eq:ch1-contraction-rate}
\pi_n = \sqrt{\frac{d}{n}}\;\vee\;n^{-\frac{\beta}{2}},
\end{equation}
其中 $\beta$ 与流形光滑度和损失函数的局部行为有关。
注意这一速率仅依赖于流形内在维度 $d$，而非环境维度 $D$——
这正是流形结构带来统计效率提升的体现。

\begin{Remark}\label[remark]{rem:ch1-rate-d}
收敛速率中 $d$ 的出现是本质的：在 $d$ 维切空间中，
Fisher 信息矩阵 $G(\theta_0)$ 的特征值个数为 $d$，
后验协方差的尺度自然地与 $d$ 成正比。
然而，与纯维数灾难不同，$d$ 是数据的有效自由度——
它是固定的（不随 $n$ 增长），因此不会导致速率的多项式衰退。
\end{Remark}

% ----------------------------------------------------------
\section{RRWM 算法}\label{sec:ch1-rrwm}

贝叶斯推断的计算核心在于从后验分布 $\Pi_n$ 中采样。
在流形约束下，标准 MCMC 算法（如随机游走 Metropolis、
Metropolis-Hastings）的直接应用面临几何障碍——
如何保证 proposal 分布产生的候选点严格位于流形上？

本文提出的 RRWM（Riemannian Random Walk Metropolis）算法
通过在切空间 $T_{\theta}\mathcal{M}$ 中构造 proposal 并利用指数映射
将其推回流形，成功解决了这一问题。

% ----------------------------------------------------------
\subsection{算法构造}\label{sec:ch1-rrwm-algo}

设当前状态为 $\theta\in\mathcal{M}$。
RRWM 的 proposal 分布定义在切空间 $T_{\theta}\mathcal{M}$ 上：
\begin{equation}\label{eq:ch1-rrwm-proposal}
v \;\sim\; N\!\big(0,\;\sigma^2\, G(\theta)^{-1}\big),
\qquad
\theta' = \exp_\theta(v),
\end{equation}
其中 $G(\theta)$ 为 $\theta$ 处的黎曼度量矩阵，
$\sigma>0$ 为步长参数。$G(\theta)^{-1}$ 作为协方差矩阵的逆
保证了 proposal 分布在黎曼几何意义下是"各向同性"的——
这被称为黎曼-高斯（ Riemannian-Gaussian）分布。

接受概率（acceptance probability）由 Metropolis 准则给出：
\begin{equation}\label{eq:ch1-rrwm-accept}
a(\theta,\theta')
= \min\!\left\{1,\;
\frac{\exp\!\big(-n L_n(\theta')\big)\,
      \pi_{\mathcal{M}}(\theta')\,q(\theta',\theta)}
     {\exp\!\big(-n L_n(\theta)\big)\,
      \pi_{\mathcal{M}}(\theta)\,q(\theta,\theta')}\right\},
\end{equation}
其中 $q(\theta,\theta')$ 为从 $\theta$ 到 $\theta'$ 的 proposal 密度
（相对于流形上的适当测度），
$\pi_{\mathcal{M}}$ 为流形体积元（volume element）。
关键的是，指数映射 $\exp_\theta$ 的 Jacobian 行列式
出现在 proposal 密度的计算中，
它补偿了流形弯曲带来的体积变化。

% ----------------------------------------------------------
\subsection{混合时间分析}\label{sec:ch1-mixing}

RRWM 算法的核心理论保证是：
其混合时间（mixing time）仅依赖流形内在维度 $d$，
而非环境维度 $D$。

\begin{Theorem}[RRWM 混合时间]\label[theorem]{thm:ch1-mixing}
在 \cref{sec:ch1-bvm} 的正则性条件下，
设 $\epsilon>0$，$n$ 固定。
RRWM 算法的 $\epsilon$-混合时间（$\epsilon$-mixing time）满足
\begin{equation}\label{eq:ch1-mixing-time}
t_{\mathrm{mix}}(\epsilon)
\;\asymp\;
d\;\cdot\;\log\!\big(d/\epsilon\big),
\end{equation}
其中隐含常数仅依赖流形几何（截面曲率、injectivity radius）
和目标分布的局部行为。
\end{Theorem}

\cref{eq:ch1-mixing-time} 的意义在于：
$D$ 完全不出现在混合时间的维度依赖中。
这与欧氏空间随机游走 Metropolis 的已知结果形成鲜明对比——
在后者的经典分析中，混合时间通常具有 $O(D)$ 的维度依赖。
这一改进来源于黎曼几何结构：
曲率约束和流形的紧性使得链在"局部"行为上不受 $D$ 的影响。

\textbf{证明思路}\footnote{详见 \citep[Section 5]{tang2023robust}。
核心思想是将流形上的链投影到一系列局部坐标卡（coordinate charts）上，
在每个局部卡内运用欧氏 Metropolis 算法的混合时间界，
然后利用流形的有限覆盖性质拼接各局部分析。
关键引理是：切空间投影的方差在 $d$ 维而非 $D$ 维尺度上增长。}：

% ----------------------------------------------------------
\section{极小极大最优性}\label{sec:ch1-minimax}

本节讨论 Bayesian RPETEL 后验的收缩速率是否达到了统计最优。
这一问题的回答需要建立上界（由 \cref{sec:ch1-rate} 的分析给出）
与相应的下界（下确界）之间的匹配。

% ----------------------------------------------------------
\subsection{收缩速率上界}\label{sec:ch1-upperbound}

由 \cref{eq:ch1-contraction-rate}，Bayesian RPETEL 后验的收缩速率上界为
$\pi_n \asymp \sqrt{d/n}$（在典型正则性假设下）。
这一速率与经典的 Le Cam's lemma 所给出的二元检验下界一致，
表明在 $L_2$ 风险意义下，$\sqrt{d/n}$ 是几乎所有估计量的最优收敛速率。

% ----------------------------------------------------------
\subsection{极小极大下界}\label{sec:ch1-lowerbound}

\begin{Theorem}[极小极大下界]\label[theorem]{thm:ch1-minimax-lower}
设 $\mathcal{M}$ 为 $d$ 维紧致光滑黎曼子流形，
族 $\mathcal{P}$ 由 $n$ 个 i.i.d. 样本的分布构成，
满足一定的光滑性假设。
则任意估计量 $\hat{\theta}_n$ 的极小极大风险满足
\begin{equation}\label{eq:ch1-minimax-lower}
\inf_{\hat{\theta}_n}\;
\sup_{P\in\mathcal{P}}\;
\mathbb{E}_P\big[\|\hat{\theta}_n-\theta^\ddagger\|_2\big]
\;\gtrsim\;
\sqrt{\frac{d}{n}}.
\end{equation}
\end{Theorem}

\cref{eq:ch1-minimax-lower} 给出了 $\sqrt{d/n}$ 的下界。
结合上界 $\pi_n \asymp \sqrt{d/n}$，我们得到：

\begin{Corollary}[极小极大最优性]\label[corollary]{cor:ch1-optimal}
在 \cref{sec:ch1-bvm} 的正则性条件下，
Bayesian RPETEL 后验的收缩速率达到了极小极大最优
（至多常数因子和对数因子）。
\end{Corollary}

\begin{Remark}\label[remark]{rem:ch1-lowerbound-tech}
\cref{eq:ch1-minimax-lower} 的证明综合运用了
Le Cam's method、tensor product construction
和黎曼几何中的 volume comparison 技术。
关键是构造一族难以区分的替代分布（alternative distributions），
其参数在流形上以适当距离分隔。
黎曼流形的体积增长性质（volume growth rate $\asymp s^d$）
在此处发挥了关键作用——它决定了在给定分辨率下
能够构造多少个相互"可区分"的备择假设。
\end{Remark}

% ----------------------------------------------------------
\section{本章小结}\label{sec:ch1-summary}

本章系统介绍了 \citep{tang2023robust} 的核心内容。
Bayesian RPETEL 框架将稳健统计的思想与黎曼几何的结构结合，
通过三重稳健性目标——似然稳健性、风险稳健性和先验稳健性——
为流形约束下的贝叶斯推断提供了统一理论。

\begin{enumerate}
    \item \textbf{RPETEL 框架}（\cref{sec:ch1-rpetel}）：
          通过伪似然构造（\cref{eq:ch1-pseudolikelihood}）和
          后验定义（\cref{eq:ch1-posterior}），
          将 PERM 嵌入贝叶斯框架，
          在平方损失下简化为对样本均值的指数型后验。
    \item \textbf{黎曼 BvM 定理}（\cref{sec:ch1-bvm}）：
          \cref{eq:ch1-bvm-limit} 建立了后验分布在切空间投影上的
          渐近正态性，是连接贝叶斯渐近理论与流形几何的核心工具。
    \item \textbf{RRWM 算法}（\cref{sec:ch1-rrwm}）：
          \cref{eq:ch1-rrwm-proposal}--\cref{eq:ch1-rrwm-accept}
          实现了流形上的 Metropolis 采样，
          \cref{eq:ch1-mixing-time} 证明其混合时间仅依赖 $d$ 而非 $D$。
    \item \textbf{极小极大最优性}（\cref{sec:ch1-minimax}）：
          \cref{cor:ch1-optimal} 表明 RPETEL 后验收缩速率
          达到极小极大最优 $\sqrt{d/n}$。
\end{enumerate}

\begin{Remark}\label[remark]{rem:ch1-connections}
本章内容与系列中其他论文存在深刻联系：
\citep{tang2022minimax} 的极小极大理论为最优性判断提供了参照系；
\citep{tang2024mala} 则利用本章建立的流形 BvM 定理，
将 MALA 在欧氏空间的最优混合性质桥接到黎曼流形 setting，
从而将后验推断的计算效率提升到了最优的 $O(d^{1/3})$ 维度依赖。
\end{Remark}
